\chapter{Conclusiones y Trabajo Futuro}

\section{Conclusiones}

El presente trabajo tuvo como objetivo principal diseñar, implementar y validar preliminarmente un sistema de microscopía computacional basado en sensor CMOS, concebido como carga útil experimental del proyecto INTISAT.

A partir de los resultados obtenidos, se concluye que:

\begin{itemize}
    \item Se implementó exitosamente un sistema de adquisición de imágenes en configuración de microscopía sin lentes, empleando un sensor CMOS comercial y una arquitectura de control programático basada en Raspberry Pi 5.

    \item Se estableció una metodología de validación metrológica preliminar mediante comparación cruzada con microscopía óptica calibrada y microscopía electrónica de barrido (SEM), asegurando una referencia dimensional independiente.

    \item El sistema demostró capacidad para detectar estructuras del orden de $2\,\mu m$, operando cerca del límite impuesto por el tamaño de píxel del sensor.

    \item Las mediciones digitales obtenidas presentan coherencia dimensional preliminar al contrastarse con referencias externas, dentro de los márgenes esperables considerando efectos de muestreo y difracción.

    \item Se identificaron explícitamente las limitaciones físicas y computacionales del sistema, incluyendo restricciones de resolución efectiva, sensibilidad a la iluminación y dependencia crítica de la distancia muestra–sensor.

    \item El procesamiento computacional mejora la interpretabilidad visual, pero no incrementa la resolución física real del sistema, lo cual fue reconocido y documentado de manera conservadora.
\end{itemize}

En conjunto, el trabajo establece una base técnica sólida para el desarrollo progresivo del payload de microscopía computacional dentro del proyecto INTISAT.
\section{Trabajo Futuro}

Los resultados obtenidos abren múltiples líneas de desarrollo orientadas a mejorar la robustez metrológica y ampliar las capacidades del sistema.

Entre las principales líneas de trabajo futuro se consideran:

\begin{itemize}
    \item Incorporación de patrones litográficos certificados (grids micrométricos) para una calibración espacial más robusta y menos sensible a efectos difractivos.

    \item Evaluación de objetos patrón de mayor tamaño (5–10 µm) que permitan validar la escala espacial en un régimen menos cercano al límite de resolución.

    \item Optimización del sistema de iluminación, incluyendo caracterización espectral y análisis cuantitativo de uniformidad mediante herramientas fotométricas.

    \item Implementación controlada de técnicas avanzadas de reconstrucción computacional (por ejemplo, FPM) bajo un marco de validación física explícita.

    \item Evaluación de sensores CMOS con menor tamaño de píxel para incrementar la resolución efectiva sin comprometer la arquitectura compacta del sistema.

    \item Desarrollo de protocolos preliminares para observación exploratoria de muestras biológicas de tamaño micrométrico (por ejemplo, estructuras celulares ampliamente utilizadas en investigación básica), sin fines diagnósticos ni clínicos.
\end{itemize}

Desde la perspectiva del proyecto INTISAT, el sistema desarrollado representa una plataforma experimental modular que puede evolucionar progresivamente hacia aplicaciones más exigentes, manteniendo coherencia metrológica y trazabilidad experimental.

El presente trabajo constituye una etapa fundacional en el desarrollo del payload, estableciendo criterios físicos, metodológicos y documentales que permitirán consolidar futuras iteraciones del sistema.
